\documentclass[10pt]{article}
\usepackage[margin=1in]{geometry}
\usepackage{amsmath,amssymb,enumitem}
\usepackage[colorlinks=true,urlcolor=blue]{hyperref}
\setlength{\emergencystretch}{1em}
\setlist[enumerate]{leftmargin=*,itemsep=0.3em}
\newcommand{\eps}{\varepsilon}
\title{Proofreading notes for SEVEM\_Stokes}
\author{}
\date{October 4, 2026}
\begin{document}
\maketitle
\noindent Scope: proof correctness, missing references, and misprints. Keep the
method, section order, and numerical data. Mark substantive manuscript changes
with \verb|\AI{...}|; leave routine wording and typo corrections unmarked.
The proof corrections and presentation edits below have been applied to
\texttt{SEVEM\_Stokes.tex} and \texttt{numerics.tex}. Author checks remain
listed separately. The original sources and PDF are saved in \texttt{proofreading/original/}.

\section*{Proof corrections}
\begin{enumerate}
\item \textbf{Mesh assumptions (Section 1.1).}
Star-shapedness alone does not make the face cones uniformly regular, and the
center of the largest inscribed ball need not be a valid star center. Choose the
center of a ball with respect to which the element is star-shaped. State the face
and cone regularity, bounded face count, and diameter comparability used by the
trace and scaling estimates. Use $h_K$ locally and $h=\max_K h_K$ globally.

\item \textbf{Velocity degrees of freedom (Section 1.2).}
Replace the set difference $\mathbb P_k(K;\mathbb R^d)\setminus\mathbf G_k(K)$
with the $L^2(K)$-orthogonal complement $\mathbf G_k(K)^\perp$ within the
polynomial vector space. The set difference is not a linear space and does not
specify an independent set of moments. State $k\geq1$ and the convention that
negative-degree polynomial spaces are zero. Cite the unisolvence result.

\item \textbf{Stress space (Section 2.1).}
The two original definitions of $\Sigma_k(T)$ differ. Extending $\mathbb P_k(F)$
constantly in the normal direction and adding $\mathbb P_{k-1}(T)$ does not give
all of $\mathbb P_k(T)$ in the normal--normal block. For $k=1$, it gives no
normal-coordinate linear term, whereas the displayed block decomposition does.
Use the block decomposition, which agrees with the stated normal bubble
$\lambda_{\boldsymbol x_K}\mathbb P_{k-1}(T)\boldsymbol n\otimes\boldsymbol n$.
This choice must be checked against the implementation.

\item \textbf{Stress unisolvence (Lemma 2.2 in the original source).}
Replace the bubble set difference by its orthogonal complement to the selected
tensor-polynomial subspace. Add the missing proof: zero face moments put the
tensor in the bubble space; zero complementary moments put it in the selected
subspace; global polynomial moments, tested with the dual tensor basis, then
force every polynomial coefficient to vanish. Count dimensions to establish
existence as well as uniqueness.

\item \textbf{Uniform stability (Assumption 2.1).}
Spanning $\mathbb S$ on each element proves algebraic independence but does not
give a mesh-independent constant. Require a uniformly conditioned choice of
normalized face-tangential tensors. Otherwise the scaling constant in the
stress inf-sup argument may diverge as the mesh changes. Prescribe tensor face
moments separately on each element side, since the stress is discontinuous.

\item \textbf{Norm and Korn arguments (Section 2.2).}
The original local equivalence with the full gradient is false: choose a
nonzero rigid rotation and matching boundary values. The strain and projected
jump vanish, but the gradient does not. Replace it by a global broken Korn
inequality, including boundary faces. In the norm proof, use both zero normal
and zero tangential boundary values; the statement that boundary normals span
$\mathbb R^d$ alone does not justify eliminating an affine rigid motion.

\item \textbf{Full versus projected jump (Sections 2.2 and 3.2).}
The norm controls $\Pi_{\partial K}\mathcal J_{\partial K}$, while the residual
contains $\mathcal J_{\partial K}$. Add
\[
h_K^{-1/2}\|\mathcal J_{\partial K}(\boldsymbol v_h)\|_{\partial K}
\lesssim\|\boldsymbol v_h\|_{1,h,K}.
\]
To prove it, subtract a local rigid motion from
$\boldsymbol w=\Pi_K^o\boldsymbol v_h^{\rm div}$, use Korn's inequality modulo
rigid motions and the trace inequality, and use the fact that the tangential
face space contains the tangential traces of rigid motions.

\item \textbf{Pressure stability and interpolation (Section 3).}
Add proofs of the pressure continuity and right-inverse lemmas. Recall the
commuting identity
\[
\operatorname{div}I_K^{\rm div}\boldsymbol v
=Q_h(\operatorname{div}\boldsymbol v)|_K.
\]
Construct the pressure lifting by interpolating an $H_0^1$ divergence lifting,
and bound its discrete norm. Since $b(\boldsymbol v_h,q_h)=
-(\operatorname{div}\boldsymbol v_h^{\rm div},q_h)$, use the negative lifting
when a positive inf-sup numerator is needed. Add absolute values to continuity
bounds and exclude zero denominators in suprema.

\item \textbf{Error equation and approximation (Section 3.2).}
Restore sums over $K$: cancellation of the exact traction on interior faces is
a global statement. The strain of $\Pi_K^o I_K^{\rm div}\boldsymbol u$ already
belongs to $\mathbb P_{k-1}(K;\mathbb S)$, so its stress projection is redundant.
Also,
\[
(\operatorname{div}\Pi_K^o\boldsymbol v_h^{\rm div},p-Q_hp)_K=0.
\]
The original estimate drops $\eps(\boldsymbol u-\Pi_K^o\boldsymbol u)$ without
justification. Retain the polynomial approximation error and bound
$\boldsymbol u-\Pi_K^o I_K^{\rm div}\boldsymbol u$ in the volume and on faces.
Estimate only the virtual interpolant's normal trace, which is a face
projection; no full $H^1$ trace of a virtual field is needed.

\item \textbf{Meaning of the final velocity error (Section 3.3).}
The norm is initially defined only on discrete pairs. Identify a smooth velocity
with $\{\boldsymbol u,\Pi_{\partial K}^t\boldsymbol u\}$ before writing
$\|\boldsymbol u-\boldsymbol u_h\|_{1,h}$. Extend the volume $L^2$ projector
to all $L^2$ vector fields. Keep the stated $O(h^k)$ estimate; do not add a claim
of pressure robustness or an $L^2$ velocity theorem.
\end{enumerate}

\section*{References to add at the point of use}
\begin{enumerate}
\item Beir\~ao da Veiga, Brezzi, Marini, and Russo,
\emph{$H(\mathrm{div})$ and $H(\mathrm{curl})$-conforming virtual element methods},
Numerische Mathematik 133 (2016), 303--332.
\href{https://doi.org/10.1007/s00211-015-0746-1}{doi:10.1007/s00211-015-0746-1}.
Use for the velocity space, moments, unisolvence, and computable projection.
\item Beir\~ao da Veiga, Mascotto, and Meng,
\emph{Interpolation and stability estimates for edge and face virtual elements
of general order}, Mathematical Models and Methods in Applied Sciences 32
(2022), 1589--1631.
\href{https://doi.org/10.1142/S0218202522500373}{doi:10.1142/S0218202522500373}.
Use as supporting interpolation and stability literature; its space variants
do not remove the need to state the estimates for the present interpolant.
\item Brenner, \emph{Korn's inequalities for piecewise $H^1$ vector fields},
Mathematics of Computation 73 (2004), 1067--1087.
\href{https://doi.org/10.1090/S0025-5718-03-01579-5}{doi:10.1090/S0025-5718-03-01579-5}.
Use for the global broken Korn inequality.
\item Boffi, Brezzi, and Fortin, \emph{Mixed Finite Element Methods and
Applications}, Springer, 2013.
\href{https://doi.org/10.1007/978-3-642-36519-5}{Publisher and DOI}.
Use for the continuous divergence lifting and mixed-method well-posedness.
\item Chen and Wang, \emph{A divergence free weak virtual element method for
the Stokes problem on polytopal meshes}, Journal of Scientific Computing 78
(2019), 864--886.
\href{https://doi.org/10.1007/s10915-018-0796-5}{doi:10.1007/s10915-018-0796-5}.
Use to identify the earlier weak VEM Stokes framework. A full comparison with stabilizer-free Stokes methods belongs in an
introduction if one is added later.
\item Mu, Ye, and Zhang, \emph{A stabilizer-free, pressure-robust, and
superconvergence weak Galerkin finite element method for the Stokes equations
on polytopal mesh}, SIAM Journal on Scientific Computing 43 (2021),
A2614--A2637.
\href{https://doi.org/10.1137/20M1380405}{doi:10.1137/20M1380405}.
Add a short reference to this related stabilizer-free Stokes method.

\end{enumerate}

\section*{Misprints and presentation}
\begin{enumerate}
\item Define the symmetric strain, rigid motions, tensor trace, inner products,
and mesh sizes before use. Distinguish a tensor's normal trace from its scalar
matrix trace. Use $K$ consistently in place of the unexplained $E$.
\item Use the existing tangential projection to shorten the weak-strain formula.
Replace the two unconditional ``$\ne$'' statements by an explanation that
integration by parts also produces internal stress-jump terms.
\item Replace the hard-coded equation number 14 and the literal ``problem (1)''
with labels. Correct inconsistent boldface, missing integration measures,
``degree of freedoms'', and repeated inequalities. Recall only the definitions
needed by a proof; keep routine language plain and short.
\item Repair undefined commands in \texttt{numerics.tex}, use a proper numerical
section heading, and give every table a label. The files
\texttt{fig/poly}, \texttt{fig/tri}, and \texttt{fig/rect} are absent. Guard the
figure input and state the missing-asset issue rather than inventing plots.
\end{enumerate}

\section*{Author checks still needed}
\begin{enumerate}
\item \textbf{Rectangular meshes are outside the stated proof.}
For an axis-aligned rectangle, the face-tangential symmetric tensors span only
the two diagonal matrices, not all of $\mathbb S$. Thus Assumption 2.1 fails.
Keep the table as an empirical result, with this limitation stated. Extending
the theory to rectangles would require additional analysis beyond this review.
The geometry of the polygonal meshes must also satisfy the strengthened
uniform assumption before the theorem is applied to them.
\item \textbf{Numerical reproducibility.}
Supply the exact velocity, pressure, force, boundary values, viscosity, mesh
construction, and definitions of the reported discrete $L^2$ and maximum norms.
The source specifies only ``sincosdata''. Confirm that the implementation uses
the analyzed strain space: the original numerical section uses
$\Sigma_0^{-1}$ while Section 2 uses $\Sigma_1$, and introduces $2\nu$ whereas
the model has coefficient one. The reported values are preserved, but their
agreement with the corrected formulation cannot be certified from these files.
\item \textbf{Front matter.}
The title command is commented out, the abstract and keywords are absent, and
the submission date and funding statement are placeholders. These require
author-supplied information; they are not filled in during proofreading.
\end{enumerate}
\end{document}
